\documentclass[11pt,letterpaper]{article}
\usepackage[T1]{fontenc}
\usepackage{lmodern}
\usepackage[left=.97in,right=.97in,top=.90in,bottom=.93in]{geometry}
\usepackage{amsmath,amssymb,amsthm,mathtools,microtype}
\usepackage{xcolor}
\usepackage[colorlinks=true,linkcolor=blue!40!black,citecolor=blue!40!black,
  urlcolor=blue!40!black]{hyperref}
\hypersetup{pdftitle={The weight-three Lambert sum at the golden ratio is irrational},
  pdfauthor={twinheterodox},pdfsubject={Number theory}}
\newtheorem{theorem}{Theorem}[section]
\newtheorem{proposition}[theorem]{Proposition}
\newtheorem{lemma}[theorem]{Lemma}
\numberwithin{equation}{section}
\newcommand{\Q}{\mathbb Q}
\newcommand{\Z}{\mathbb Z}
\newcommand{\C}{\mathbb C}
\newcommand{\OK}{\mathcal O}
\DeclareMathOperator{\Nm}{N}
\title{\fontsize{16}{20}\selectfont
  The weight-three Lambert sum at the golden ratio is irrational}
\author{twinheterodox}
\date{18 November 2026}

\begin{document}
\maketitle

\begin{abstract}
For \(|b|>1\), define the weight-three Lambert sum by
\(Z(b)=\sum_{n\ge1}n^2/(b^n-1)\). For an algebraic integer \(b\)
in this range, we prove that \(Z(b)\notin\Q(b)\) if every conjugate
outside the closed unit disk belongs to \(\{b,\bar b\}\). This includes
all Pisot and Salem bases. For a Salem base we bound the polynomial
approximations even at conjugates on the unit circle, where the series
diverges. At the golden ratio \(\varphi=(1+\sqrt5)/2\), the sum
\(Z(\varphi)\) lies outside \(\Q(\sqrt5)\) and is therefore irrational.
A nonreal example is \(Z(1+i)\notin\Q(i)\).
\end{abstract}

\section{The golden-ratio value}\label{sec:intro}

For a complex number \(b\) with \(|b|>1\), define
\begin{equation}\label{eq:lambert}
 Z(b)=\sum_{n=1}^{\infty}\frac{n^2}{b^n-1}
     =\sum_{m=1}^{\infty}\frac{\sigma_2(m)}{b^m},
 \qquad \sigma_2(m)=\sum_{d\mid m}d^2 .
\end{equation}
Expanding each denominator as a geometric series gives the second
expression. Both series converge absolutely. At the golden ratio
every partial sum belongs to the quadratic field \(\Q(\sqrt5)\).

\begin{theorem}\label{thm:gold}
For \(\varphi=(1+\sqrt5)/2\),
\[
                     \sum_{n=1}^{\infty}\frac{n^2}{\varphi^n-1}
                         \notin \Q(\sqrt5).
\]
In particular this real number is irrational.
\end{theorem}

The second root of \(x^2-x-1\) is \(-1/\varphi\). Suppose the Lambert
sum at \(\varphi\) were \(\xi=u+v\sqrt5\) with \(u,v\in\Q\).
Its conjugate would be \(Y=u-v\sqrt5\). Substituting \(-1/\varphi\)
in the Lambert series makes its \(n\)th term asymptotic to \(-n^2\).
To compare the two embeddings, we construct integer polynomials
\(P_K(t,X)\) of degree at most \(K\) in \(X\). For some \(c>0\) we
have \(|P_K(\varphi,\xi)|\le\exp(-cK^3)\) for all large \(K\). At
\((-1/\varphi,Y)\) we bound their growth by \(\exp(o(K^3))\)
\emph{for every fixed \(Y\)}. Choose a positive integer \(d\) for
which \(d\xi\) is an algebraic integer. Then
\(d^{2K}P_K(\varphi,\xi)P_K(-1/\varphi,Y)\) is an integer. Its absolute value
tends to zero, and we show it is nonzero for arbitrarily large \(K\).

For a nonreal base \(b\), absolute convergence gives
\(Z(\bar b)=\overline{Z(b)}\), so decay at \(b\) gives decay at
\(\bar b\) as well. It remains to control the other conjugates, which
we can do whenever they all lie in the closed unit disk.

\begin{theorem}\label{thm:main}
Let \(b\) be an algebraic integer with \(|b|>1\). If every conjugate
distinct from \(b\) and its complex conjugate \(\bar b\) has modulus
at most one, then \(Z(b)\notin\Q(b)\).
\end{theorem}

A real algebraic integer \(b>1\) is a Pisot number when all its
other conjugates have modulus less than one, and a Salem number when
they lie in the closed unit disk with at least one on the unit circle.
Theorem~\ref{thm:main} covers both classes. At a unit-circle conjugate
\(u\) of a Salem number, \(|u^n-1|\le2\), so the terms of the series
cannot tend to zero. Section~\ref{sec:disk} controls the polynomials there.

The Gaussian integer \(1+i\) has norm \(2\), the smallest norm above
one in \(\Z[i]\). Its only other conjugate is \(1-i\), so
\(Z(1+i)\notin\Q(i)\). Every partial sum has rational real and imaginary
parts. At least one of those coordinates is irrational in the limit.

In the Lambert-series convention for \(q\)-zeta values,
\(Z(b)=\zeta_{1/b}(3)\). Bundschuh and Zudilin singled out reciprocal
Pisot and Salem parameters in their discussion of the arithmetic of
these values \cite[\S6, after (41)]{BZ}.
For the weight-one sum \(\sum_{n\ge1}1/(b^n-1)\), Tachiya's
Theorems~1 and~2 give exclusion from \(\Q(b)\) at \(b=1+i\) and
\(b=\varphi\), respectively \cite[p.~232]{Tachiya}. Dixit, Kumar and
Pathak prove this exclusion at every Pisot base \cite[Corollary~1.2]{DKP}.
For the weight-three sum \eqref{eq:lambert}, uplift-universal-inc
proved irrationality at each positive integer base \(b\ge2\)
\cite[Theorem~1.1]{Uplift}.

The determinant and its cyclotomic factors come from the integer-base
proof \cite[\S\S1--5 and Appendix~A]{Uplift}. Zudilin obtained cubic
decay for Hankel determinants on geometric nodes in his work on the
generalized \(q\)-logarithm \cite[\S4, (8)]{Zudilin}.
Krisman used a norm over a quadratic field and controlled the conjugate
of a determinant for reciprocal Fibonacci cubes \cite[\S5]{Krisman}.
Here a single estimate controls \(P_K(t,Y)\) throughout the closed
unit disk for every fixed \(Y\). We estimate the determinant on a
slightly larger circle, where successive powers of \(t\) are
separated in absolute value. Since \(P_K\) is a polynomial in \(t\), the maximum
modulus principle carries that estimate to every remaining conjugate.

\section{A single integral polynomial for every embedding}\label{sec:poly}

Put \(q=b^{-1}\). To recover the summand \(j^2/(b^j-1)\) from
\(x^j\), apply the Euler derivative \(\theta=x\,d/dx\) twice and sum
over the points \(q,q^2,\ldots\). Thus for a function \(f\)
holomorphic on a neighborhood of \(\lvert x\rvert\le |q|\), set
\[
                  L_q(f)=\sum_{m\ge1}(\theta^2 f)(q^m).
\]
Since \(\theta^2 f(x)=xf'(x)+x^2 f''(x)\), the summands are
\(O_f(|q|^m)\). On monomials this gives
\[
 L_q(1)=0,\qquad L_q(x^j)=\frac{j^2}{b^j-1}\quad(j\ge1).
\]
The same functional recovers \emph{tails} of \(Z(b)\) on simple
fractions. Indeed, with
\[
 g(y)=\sum_{j\ge1}j^2 y^j=\frac{y(1+y)}{(1-y)^3}\qquad(|y|<1),
\]
absolute summation gives \(Z(b)=\sum_{m\ge1}g(q^m)\) and
\begin{equation}\label{eq:tail}
 L_q\big((1-q^a x)^{-1}\big)=\sum_{m\ge1}g(q^{m+a})
    =Z(b)-S_a(b),\qquad
 S_a(t)=\sum_{m=1}^a\frac{t^m(t^m+1)}{(t^m-1)^3}.
\end{equation}
For example the first pole, at \(x=b\), removes just \(g(q)\)
from \(Z(b)\). Thus every pole introduces the same unknown value.

Now replace the base \(b\) and its value \(Z(b)\) by independent
variables \(t\) and \(X\), respectively. The functional still acts
in \(x\). For \(K\ge1\) the first \(K\) poles give the denominator
\[
                    D_{K,t}(x)=\prod_{a=1}^K(1-t^{-a}x).
\]
On functions \(p(x)/D_{K,t}(x)\), define \(\mathcal L_{t,X}\) by
partial fractions over \(\Q(t,X)\). On the polynomial part use
\(\mathcal L_{t,X}(1)=0\) and
\(\mathcal L_{t,X}(x^j)=\mu_j(t)=j^2/(t^j-1)\) for \(j\ge1\).
On the simple fractions use
\(\mathcal L_{t,X}((1-t^{-a}x)^{-1})=X-S_a(t)\).
By \eqref{eq:tail}, at \(t=b,\ X=Z(b)\) this agrees with \(L_q\).

\begin{samepage}
We arrange the moments in a determinant so that we can change the
polynomial basis in both its rows and its columns. At \(b\), monic
polynomials with triple zeros at the first few nodes \(q,q^2,\ldots\)
remove those nodes even after two derivatives. Successive rows and
columns then supply the small factors that give cubic decay.
In the monomial basis the determinant is
\begin{equation}\label{eq:Delta}
 \Delta_{K,t}(X)=\det\left[
    \mathcal L_{t,X}\big(x^{i+j}/D_{K,t}(x)\big)
                         \right]_{0\le i,j<K}.
\end{equation}
\end{samepage}
The entries depend only on \(i+j\), giving a Hankel matrix. Each is
affine in \(X\), so the determinant has degree at most \(K\) in \(X\).
Rational modifications of moment functionals appear in Uvarov's
transformation of orthogonal polynomials \cite[\S1]{Uvarov} and in
Krattenthaler's determinant identity \cite[Theorem~1]{Kratt}.
The denominator \(D_{K,t}\) and the functional in \eqref{eq:Delta}
are the specific weight-three construction of \cite[\S1]{Uplift}.

Two features of this construction are essential for specialization.
The poles of the coefficients in \(t\), away from zero, occur at
roots of unity where nodes coincide or moments have poles.
Cyclotomic factors clear them, and \(t^{e_K}\) divides all the
resulting coefficients, where \(e_K=(K^3-K)/2\).
Let \(\Phi_d\) be the \(d\)-th cyclotomic
polynomial. In units of \(K^3\log|b|\), the determinant will save
\(1/9\) and the extracted power of \(t\) another \(1/2\).
The cost of clearing the denominators comes from the degree of
\begin{equation}\label{eq:clearer}
 C_K(t)=\prod_{d=1}^K \Phi_d(t)^{c_{K,d}},\qquad
 c_{K,d}=
 \begin{cases}
  \left\lceil K^2/(4d)+2K+17d/16\right\rceil,& d\le K/2,\\
  5(K-d)+3,& K/2<d\le K .
 \end{cases}
\end{equation}

\begin{lemma}\label{lem:arithmetic}
The polynomial used at every embedding is
\[
             P_K(t,X)=C_K(t)t^{-e_K}\Delta_{K,t}(X)\in\Z[t,X].
\]
It has degree at most \(K\) in \(X\). Moreover,
\begin{equation}\label{eq:degree}
 \deg C_K\le\gamma K^3+O(K^2\log(2K)),\qquad
 \gamma=\frac{321}{64\pi^2},\qquad
 \sum_{d=1}^K c_{K,d}=O(K^2\log(2K)).
\end{equation}
\end{lemma}

\begin{proof}
At a primitive \(d\)-th root of unity, \(t^a=t^c\) precisely when
\(a\equiv c\pmod d\). The first bound in \eqref{eq:clearer} comes
from \cite[\S4.2, (4.8)]{Uplift}. For \(d>K/2\) each class contains
at most two indices, and \cite[Proposition~A.1]{Uplift} gives the
second bound. Combining these local bounds as in \cite[\S4.4]{Uplift} gives
\(C_K(t)\Delta_{K,t}(X)\in\Z[t,t^{-1}][X]\).
Their cyclotomic degree sum gives the first bound in \eqref{eq:degree}
\cite[\S4.5 and Proposition~A.1]{Uplift}. The range \(d\le K/2\)
contributes \(161/(64\pi^2)\) to its leading coefficient and the
range \(d>K/2\) contributes \(5/(2\pi^2)\). The bound on the sum of
exponents follows directly from \eqref{eq:clearer}.

It remains to justify using the extracted factor \(t^{e_K}\) as an
identity in the base. We already know that the coefficients of \(P_K\)
are integer Laurent polynomials. The integer-base divisibility result
\cite[Proposition~3.1 and Appendix~A]{Uplift} gives
\(P_K(p,X)\in\Z[X]\) for every positive rational prime \(p\).
If a coefficient had a lowest term \(c t^{-h}\) with \(h>0\) and
\(c\ne0\), choose a prime \(p\nmid c\). Multiplying its value at
\(p\) by \(p^h\) gives an integer congruent to \(c\bmod p\), which
contradicts the integrality of the value itself. This proves the
assertion over \(\Z[t,X]\). The inverse-base polynomial normalization
is also recorded in the rational-parameter companion to
\cite[\S7.1]{Uplift}. The code in \cite{code} checks this integrality
coefficientwise for \(1\le K\le4\).
\end{proof}

For later use, if \(t\ne0\) is fixed and \(|t|\ne1\), Möbius
inversion in the cyclotomic product gives
\begin{equation}\label{eq:cyclotomic-away}
  \log|\Phi_d(t)|=\phi(d)\log^+|t|+O_t(1),\qquad
  \log^+u=\log\max(1,u).
\end{equation}
Here \(\phi\) is Euler's totient. After factoring off the degree
term, the absolute error is at most
\(\sum_{h\ge1}-\log(1-r^h)\) with \(r=\min(|t|,|t|^{-1})<1\).
The bound is independent of \(d\).

\section{Decay where the Lambert series converges}\label{sec:expanding}

Fix any complex \(b\) of modulus greater than one and put \(q=1/b\),
\(r=|q|\). At \(X=Z(b)\), the determinant in \eqref{eq:Delta} is
formed from the bilinear form \(L_q(fg/D_{K,b})\). Its first evaluation
node is \(x=q\). A factor \((x-q)^3\) makes the value and the first
two derivatives vanish at that node. Three more factors can remove
the node \(q^2\), and so on. This accounts for the multiplicity
three in the following basis from \cite[\S2]{Uplift}.

\begin{proposition}\label{prop:decay}
For every \(|b|>1\) and every \(K\ge1\),
\[
 |P_K(b,Z(b))|\le
 \exp\bigl(-\varepsilon K^3\log|b|+O_b(K^2\log(2K))\bigr),
 \qquad \varepsilon=\frac{11}{18}-\frac{321}{64\pi^2}>0.
\]
\end{proposition}

\begin{proof}
Define \(B_0=1\) and, for \(j\ge1\),
\[
 B_j(x)=\prod_{\ell=1}^j(x-q^{\lceil\ell/3\rceil}),\qquad
 s_j=\sum_{\ell=1}^j\lceil\ell/3\rceil,\qquad R_j=\lceil j/3\rceil.
\]
Set \(s_0=R_0=0\). For example, \(B_6=(x-q)^3(x-q^2)^3\) and
\(s_6=3\cdot1+3\cdot2=9\). More generally,
\(s_{3h}=3(1+\cdots+h)\). Thus the exponent contributed by a row
grows quadratically with its degree. Adding the exponents from all
rows and columns will give the cubic power of \(K\).

The change from monomials to the \(B_j\) is
triangular with ones on its diagonal, so it preserves
\(\Delta_{K,b}(Z(b))\). In the entry indexed by \(i,j\), put
\(R=\max(R_i,R_j)\) and \(W=1/D_{K,b}\). For \(1\le m<R\)
the node \(q^m\) contributes zero to \(L_q(B_iB_jW)\) because it
is a triple zero of \(B_iB_j\).

For \(m\ge\max(1,R)\), each factor \(x-q^a\) at \(x=q^m\) has
modulus at most \(2r^a\). This supplies \(r^{s_i+s_j}\) in the
product. To sum over the remaining nodes we also need decay in \(m\).
Both \(\theta(x-q^a)\) and \(\theta^2(x-q^a)\) equal \(x\).
At \(q^m\) their modulus is at most \(r^a r^{m-R}\), since \(a\le R\).
On the unit disk the function \(W\) is bounded
by \(\prod_{a\ge1}(1-r^a)^{-1}\). Cauchy's estimates then give, on
the smaller disk \(|x|\le r\),
\[
       |W(x)|\le C_r,\qquad
       |\theta W(x)|+|\theta^2W(x)|\le C_r |x|.
\]
Consequently every term of \(\theta^2(B_iB_jW)(q^m)\) has the
extra summable factor \(r^{m-R}\), including the term in which
both derivatives fall on \(W\). There are at most \((i+j+1)^2\)
ordered choices of the differentiated factors. Summing over \(m\)
gives
\[
       |L_q(B_iB_jW)|\le C'_r(i+j+1)^2\,2^{i+j}r^{s_i+s_j}.
\]
Take \(2^i r^{s_i}\) from each row and \(2^j r^{s_j}\) from each
column. The remaining determinant has absolute value at most
\(K![C'_r(2K)^2]^K\). The extracted power of \(r\) has exponent
\(2\sum_{j<K}s_j=K^3/9+O(K^2)\), since \(s_j=j^2/6+O(j)\). Thus
\begin{equation}\label{eq:rawdecay}
    |\Delta_{K,b}(Z(b))|
       \le \exp\bigl(-\tfrac19 K^3\log|b|+O_b(K^2)\bigr).
\end{equation}
This estimate uses only the moduli of the geometric nodes.
Multiplying by \(b^{-e_K}\) saves another \((K^3-K)\log|b|/2\)
in the logarithm. Equations \eqref{eq:degree} and
\eqref{eq:cyclotomic-away} bound the logarithm of the cost of
\(C_K(b)\) by \(\gamma K^3\log|b|+O_b(K^2\log(2K))\).
The claimed margin is \(1/9+1/2-\gamma>0\), since \(\pi>3\).
\end{proof}

\section{One bound on the closed unit disk}\label{sec:disk}

At a conjugate such as \(-1/\varphi\), the norm argument needs only
an upper bound \(o(K^3)\) for \(\log^+|P_K(t,Y)|\) at each fixed \(Y\).
The individual moments can have poles in the
closed unit disk. For example, \(S_1(t)\) has a pole at \(t=1\), yet
\(P_1(t,Y)=(t-1)^3Y-t(t+1)\) gives \(P_1(1,Y)=-2\).
Since \(P_K(\cdot,Y)\) is a polynomial, we can estimate it on a
circle just outside the disk.

\begin{lemma}\label{lem:disk}
For every fixed \(Y\in\C\),
\[
   \max_{|t|\le1}\log^+|P_K(t,Y)|
        = O_Y\!\left(K^2\log^3(2K)\right).
\]
\end{lemma}

\begin{proof}
Choose \(|t|=1+1/K\). The outward shift keeps the nodes
\(r_a=t^a\), \(1\le a\le K\), bounded by \(e\). By \eqref{eq:degree},
the degree term \((\deg C_K)\log|t|\) costs only \(O(K^2)\) at this
radius. At the same time,
distinct nodes are separated uniformly in the argument of \(t\). For
\(a\ne c\) the reverse triangle inequality gives
\[
    |1-r_a/r_c|\ge \big|1-|t|^{a-c}\big|\ge \frac1{K+1}.
\]
Consequently the weights in the partial fractions of \(1/D_{K,t}\)
satisfy
\[
 \frac1{D_{K,t}(x)}=\sum_{a=1}^K\frac{w_a}{1-x/r_a},
 \qquad w_a=\prod_{c\ne a}(1-r_a/r_c)^{-1},
 \qquad |w_a|\le(K+1)^{K-1}.
\]

To bound a moment with numerator \(x^k\), put
\(U_{a,k}=\mathcal L_{t,Y}(x^k/(1-x/r_a))\) and set \(\mu_0=0\).
The identity \(x/(1-x/r_a)=r_a((1-x/r_a)^{-1}-1)\) gives
\[
 U_{a,0}=Y-S_a(t),\qquad
 U_{a,k+1}=r_a\bigl(U_{a,k}-\mu_k(t)\bigr).
\]
For \(j\ge1\) we have \(|t^j-1|\ge1/K\). In \eqref{eq:tail}
there are at most \(K\) terms, each of size \(O(K^3)\), since
\(|t^m|\le e\) for \(m\le K\). Hence \(|S_a(t)|=O(K^4)\), and
\(|\mu_j(t)|\le4K^3\) for \(1\le j\le2K\).
Iterating the recurrence yields
\(|U_{a,k}|\le e^{2K}(|Y|+O(K^4))\) throughout \(0\le k\le2K-2\).
Each entry in \eqref{eq:Delta} is \(\sum_a w_aU_{a,k}\) with
\(k=i+j\), so its absolute value is at most \(\exp(O_Y(K\log(2K)))\).
Expanding the determinant gives
\[
                \log^+|\Delta_{K,t}(Y)|=O_Y(K^2\log(2K)).
\]

It remains to bound the cyclotomic product on the same circle.
For \(h\ge1\), we have \(1/(K+1)\le |1-t^{-h}|\le2\).
Möbius inversion for \(\Phi_d\) therefore gives
\[
 \log|\Phi_d(t)|\le
    \phi(d)\log(1+1/K)+\tau(d)\log(2K+2),
\]
where \(\tau(d)\) is the number of positive divisors of \(d\).
After weighting by \(c_{K,d}\), the first terms sum to the degree
contribution above. For the divisor terms, \eqref{eq:clearer}
gives \(c_{K,d}\le8K^2/d\)
and
\[
 \sum_{d\le K}\frac{\tau(d)}d
   =\sum_{uv\le K}\frac1{uv}\le (1+\log K)^2.
\]
Thus \(\log^+|C_K(t)|=O(K^2\log^3(2K))\). Since
\(|t^{-e_K}|\le1\), this also bounds \(\log^+|P_K(t,Y)|\)
on the circle, with an implied constant depending on \(Y\).
Lemma~\ref{lem:arithmetic} makes \(t\mapsto P_K(t,Y)\) a polynomial.
The maximum modulus principle gives the same bound on the closed unit
disk, including the points where the original moment formulas have poles.
\end{proof}

\section{Infinitely many nonzero determinants}\label{sec:nonzero}

To use the field norm we need arbitrarily large sizes at which the
determinant is nonzero. For the functional \(L_q\), which can take
complex values, we detect nonvanishing through its Cauchy transform.
For a fixed moment sequence, Kronecker related eventual vanishing of
its Hankel determinants to rationality of its generating function
\cite[Arts.~VII and~IX]{Kronecker}. Zudilin used this criterion for
the generalized \(q\)-logarithm \cite[\S4]{Zudilin}. Here the
denominator changes with the size of the determinant.

The zeros \(b,b^2,\ldots,b^K\) of \(D_{K,b}\) form nested interpolation
sets. If all sufficiently large sizes were singular, one rational
function would agree with the Cauchy transform at infinitely many of
these nodes. The nodes tend to infinity, so the two functions would
have the same Laurent expansion there. This is impossible because the
transform has a pole at every \(b^{-m}\).

\begin{proposition}\label{prop:nonzero}
For every complex \(b\) with \(|b|>1\), there are arbitrarily large
\(K\) for which \(P_K(b,Z(b))\ne0\).
\end{proposition}

\begin{proof}
We adapt \cite[Proposition~5.1]{Uplift} to complex nodes. Put
\(q=1/b\) and let
\[
 \mathcal C(z)=L_q\big((z-x)^{-1}\big)
       =\sum_{m\ge1}\frac{q^m(z+q^m)}{(z-q^m)^3}
                  \qquad(|z|>|q|).
\]
On compact sets disjoint from zero and the nodes, the tail is
dominated by a constant times \(\sum_m|q|^m\), so the series
continues meromorphically to the punctured plane. The nodes have
distinct moduli, and at each \(q^m\) the coefficient of
\((z-q^m)^{-3}\) is \(2q^{2m}\ne0\). Thus the function analytic
near infinity has infinitely many poles on continuation and cannot
be rational.

Suppose the determinant of size \(k\) is zero. Write \(D_k=D_{k,b}\).
A nonzero null vector supplies a polynomial \(Q_k\) of degree
less than \(k\) satisfying
\[
           L_q(x^jQ_k(x)/D_k(x))=0\qquad(0\le j<k).
\]
We seek a polynomial \(T_k\) such that the error
\(\mathcal C Q_k-T_k\) vanishes at \(b,\ldots,b^k\) and tends to
zero at infinity. To impose the zeros, multiply the Cauchy transform
of \(Q_k(x)/D_k(x)\) by \(D_k(z)\). Subtracting this error from
\(\mathcal C(z)Q_k(z)\) gives \(T_k\). With \(L_q\) acting in \(x\),
the two expressions are
\begin{align}
 E_k(z)&=D_k(z)L_q\!\left(\frac{Q_k(x)}{(z-x)D_k(x)}\right),
                                                     \label{eq:error}\\
 T_k(z)&=\mathcal C(z)Q_k(z)-E_k(z)
       =L_q\!\left(\frac{Q_k(z)D_k(x)-D_k(z)Q_k(x)}
                               {(z-x)D_k(x)}\right).\notag
\end{align}
The numerator defining \(T_k\) is divisible by \(z-x\), so \(T_k\)
is a polynomial of degree less than \(k\). The functional in
\eqref{eq:error} is holomorphic at every \(z=b^a,\ 1\le a\le k\),
giving the desired zeros. At infinity expand \(1/(z-x)\) to \(k\)
terms with remainder
\(x^k/(z^k(z-x))\). All \(k\) terms disappear by the moment
equations. To bound the remainder, choose a disk about zero of
radius strictly between \(|q|\) and one. Throughout this disk
the factor \(x^k Q_k(x)/((z-x)D_k(x))\) is \(O_k(|z|^{-1})\).
Cauchy's inequalities bound its first two derivatives on
\(|x|\le |q|\), and the Euler derivatives supply a summable factor
\(|q|^m\) at \(q^m\). The remainder therefore contributes
\(O_k(|z|^{-k-1})\) to the functional, giving \(E_k(z)=O_k(z^{-1})\)
in \eqref{eq:error}. The same local bounds in \(z\) justify
holomorphy at the exterior zeros.

If sizes \(k\) and \(k+1\) both vanish, the polynomial
\[
      T_{k+1}Q_k-T_kQ_{k+1}=E_kQ_{k+1}-E_{k+1}Q_k
\]
has degree at most \(k-1\), by its growth at infinity, and vanishes
at the \(k\) common exterior nodes. It is identically zero.
Hence if all sufficiently large sizes vanished, the fractions
\(T_k/Q_k\) would all represent one rational function \(T/Q\)
in lowest terms, allowing \(T=0,\ Q=1\) for the zero function.
Then \(T_k=V_kT,\ Q_k=V_kQ\) for a nonzero polynomial \(V_k\).
If \(V_k(b^a)=0\), the interpolation condition \(E_k(b^a)=0\) holds
whatever value \(\mathcal C\) takes there. We must bound how many
interpolation nodes this common factor can account for.

Since \(\mathcal C\) is nonrational, \(E=\mathcal C Q-T\) is meromorphic
at infinity and has a nonzero leading term there. The estimate
\(V_kE=O_k(1/z)\) forces
\(E(z)\sim c z^{-h}\) there for some \(c\ne0\) and integer \(h\ge1\).
Thus \(\deg V_k\le h-1\) for every sufficiently large \(k\), and
\(E(b^a)\ne0\) for every sufficiently large \(a\).
Then \eqref{eq:error} forces \(V_k(b^a)=0\) for all these \(a\le k\).
The number of distinct zeros grows beyond \(h-1\). This contradiction
proves that \(\Delta_{K,b}(Z(b))\ne0\) for arbitrarily large \(K\).
The factors relating \(P_K\) to \(\Delta_{K,b}\) are nonzero at \(b\).
\end{proof}

\section{The field norm}\label{sec:norm}

\begin{proof}[Proof of Theorem~\ref{thm:main}]
Suppose that \(\xi=Z(b)\) belongs to \(E=\Q(b)\). Choose a
positive integer \(d\) such that \(d\xi\) is in the ring of
integers \(\OK_E\). Since \(b\in\OK_E\) and
\(P_K\in\Z[t,X]\) has \(X\)-degree at most \(K\), we have
\[
                       y_K=d^K P_K(b,\xi)\in\OK_E .
\]
By Proposition~\ref{prop:nonzero}, \(y_K\ne0\) along an unbounded
sequence of sizes. Its absolute norm at each such size is a positive
integer, given by
\[
  |\Nm_{E/\Q}(y_K)|
    =d^{K[E:\Q]}\prod_{\sigma:E\hookrightarrow\C}
                  |P_K(\sigma b,\sigma\xi)|.
\]
At the chosen embedding, Proposition~\ref{prop:decay} supplies a
negative term \(-\varepsilon K^3\log|b|\) in the logarithmic upper
bound. When \(b\) is nonreal, the restriction of complex conjugation
to \(E\) is the embedding into \(\C\) that sends \(b\) to \(\bar b\). It sends
\(\xi\) to \(\overline{Z(b)}=Z(\bar b)\) and gives the same bound.
For every other embedding, \(|\sigma b|\le1\). Lemma~\ref{lem:disk},
with the fixed argument \(Y=\sigma\xi\), bounds the total logarithmic
cost of these finitely many factors by \(O_{b,\xi}(K^2\log^3(2K))=o(K^3)\).

Let \(r_b=1\) for real \(b\) and \(r_b=2\) otherwise. The fixed
denominator contributes only \(K[E:\Q]\log d\). On the nonzero
subsequence we obtain
\[
   0\le \log|\Nm_{E/\Q}(y_K)|
       \le -r_b\varepsilon(\log|b|)K^3+o(K^3),
\]
which is impossible for large \(K\). This proves
Theorem~\ref{thm:main} and, at \(b=\varphi\),
Theorem~\ref{thm:gold}.
\end{proof}

\section*{Verification}
Another AI model independently audited the mathematics and exposition of
the initial 18 November 2026 manuscript and found no gap in either
theorem. The uniform disk estimate in Section~\ref{sec:disk} follows
the proof supplied in that audit.

\small
\begin{thebibliography}{10}
\setlength{\itemsep}{1pt}
\setlength{\parskip}{0pt}
\bibitem{BZ}
P.~Bundschuh and W.~Zudilin,
Irrationality measures for certain \(q\)-mathematical constants,
\emph{Math. Scand.} \textbf{101} (2007), 104--122.
\href{https://doi.org/10.7146/math.scand.a-15034}{doi:10.7146/math.scand.a-15034}.

\bibitem{DKP}
A.~B.~Dixit, V.~Kumar and S.~S.~Pathak,
Linear independence of values of the \(q\)-exponential and related
functions, \emph{Bull. Aust. Math. Soc.} \textbf{109} (2024), 453--463.
\href{https://doi.org/10.1017/S0004972723001028}{doi:10.1017/S0004972723001028}.

\bibitem{Kratt}
C.~Krattenthaler,
A determinant identity for moments of orthogonal polynomials that
implies Uvarov's formula for the orthogonal polynomials of rationally
related densities, \href{https://arxiv.org/abs/2103.03969}{arXiv:2103.03969}
(2021).

\bibitem{Krisman}
S.~Krisman,
The sum of the reciprocal cubes of the Fibonacci numbers is irrational,
\href{https://hexagonmath.org/2611.00473v1}{Hexagon 2611.00473v1} (2026).

\bibitem{Kronecker}
L.~Kronecker, Zur Theorie der Elimination einer Variabeln aus zwei
algebraischen Gleichungen, \emph{Monatsber. K\"onigl. Preuss.
Akad. Wiss. Berlin} (1881), 535--600.
\href{https://www.e-rara.ch/download/pdf/5761162.pdf\#page=42}{Art.~IX in the collected works}.

\bibitem{Tachiya}
Y.~Tachiya, Irrationality of certain Lambert series,
\emph{RIMS K\^oky\^uroku} \textbf{1384} (2004), 231--238.
\href{https://www.kurims.kyoto-u.ac.jp/~kyodo/kokyuroku/contents/pdf/1384-34.pdf}{Kyoto account}.
See also \emph{Tokyo J. Math.} \textbf{27} (2004), 75--85,
\href{https://doi.org/10.3836/tjm/1244208475}{doi:10.3836/tjm/1244208475}.

\bibitem{Uplift}
uplift-universal-inc, Irrationality of the Lambert-series value
\(\zeta_{1/p}(3)\), research draft, 12 October 2026.
\href{https://github.com/uplift-universal-inc/qzeta_3-irrationality/blob/main/paper/main.pdf}{Manuscript},
\href{https://github.com/uplift-universal-inc/qzeta_3-irrationality/blob/main/docs/FUTURE_RESULTS.md}{companion results}.

\bibitem{Uvarov}
V.~B.~Uvarov, The connection between systems of polynomials orthogonal
with respect to different distribution functions,
\emph{USSR Comput. Math. Math. Phys.} \textbf{9} (1969), no.~6, 25--36.
\href{https://doi.org/10.1016/0041-5553(69)90124-4}{doi:10.1016/0041-5553(69)90124-4}.

\bibitem{Zudilin}
W.~Zudilin, On the irrationality of generalized \(q\)-logarithm,
\emph{Res. Number Theory} \textbf{2} (2016), article~15.
\href{https://doi.org/10.1007/s40993-016-0042-x}{doi:10.1007/s40993-016-0042-x}.

\bibitem{code}
twinheterodox, \href{https://github.com/twinheterodox/algebraic-base-lambert-values}{Verification code for the weight-three Lambert sum}, GitHub (2026).
\end{thebibliography}
\end{document}
